\documentclass[../../../include/open-logic-section]{subfiles}

\begin{document}

\olfileid[tr]{sth}{story}{predicative}

\olsection{Yüklemsel ve Yüklemsel Olmayan}

Russell kümesi $R$, $\Setabs{x}{x\notin x}$ aracılığıyla tanımlandı. Daha
açık söylersek $R$, tam olarak kendisinin elemanı olmayan bütün kümeleri
içeren küme olacaktı. Dolayısıyla $R$'yi tanımlarken (var olsaydı) $R$'yi de
içerecek alan üzerinde niceleme yaparız.

Bu \emph{yüklemsel olmayan} bir tanımdır. Daha genel olarak, bir tanım ancak
ve ancak tanımlanan nesneyi içeren bir alan üzerinde niceleme yapıyorsa
yüklemsel değildir diyebiliriz.

Paradoksların ardından Whitehead, Russell, Poincar\'{e} ve Weyl bu tür
yüklemsel olmayan tanımları ``kısır döngülü'' bularak reddetti:
\begin{quote}
  Kaçınılması gereken paradoksların çözümlemesi, hepsinin bir tür kısır
  döngüden kaynaklandığını gösterir. Söz konusu kısır döngüler, bir nesneler
  koleksiyonunun ancak koleksiyonun bütünü aracılığıyla tanımlanabilen
  elemanlar içerebileceğini varsaymaktan doğar[\ldots. \textparagraph]

  Gayrimeşru bütünlüklerden kaçınmamızı sağlayan ilke şöyle ifade edilebilir:
  ``Bir koleksiyonun \emph{tümünü} içeren şey o koleksiyonun üyelerinden biri
  olmamalıdır''; ya da tersine: ``Belli bir koleksiyonun bir bütünü olsaydı,
  yalnızca o bütün aracılığıyla tanımlanabilen üyeler içerirdi denebiliyorsa,
  söz konusu koleksiyonun bütünü yoktur.'' Buna ``kısır döngü ilkesi''
  diyeceğiz; çünkü gayrimeşru bütünlük varsayımının içerdiği kısır
  döngülerden kaçınmamızı sağlar. \citep[p.~37]{WhiteheadRussell1910}
\end{quote}
Onları izleyip \emph{kısır döngü ilkesini} benimsersek,
\olref[sth][story][rus]{sec}'deki yıkıcı Naif Küme Oluşturma Şeması'nı
şuna benzer bir şeyle değiştirmeyi deneyebiliriz:

\begin{defish}
\emph{Yüklemsel Küme Oluşturma.} Yalnızca kümeler üzerinde niceleme yapan
her $\phi$ formülü için küme$^\prime$ $\Setabs{x}{\phi(x)}$ vardır.
\end{defish}

Kümeler$^{\prime}$ küme olmadığı sürece çelişki doğmaz.

Ne yazık ki Yüklemsel Küme Oluşturma pek \emph{kapsamlı} değildir. Bizi yeni
varlıklarla, kümeler$^\prime$ ile tanıştırır. Öyleyse kümeler$^\prime$
üzerinde niceleme yapan formülleri de ele almamız gerekir. Bunlar her zaman
bir küme$^\prime$ verirse, kendisinin elemanı olmayan bütün
kümeler$^\prime$ın küme$^\prime$ını düşünerek Russell paradoksunu yeniden
elde ederiz. Aynı düşünceyi sürdürürsek kümeler$^\prime$ üzerinde niceleme
yapan bir formülün karşılık gelen bir küme$^{\prime\prime}$ verdiğini
söylemeliyiz. Sonra kümeler$^{\prime\prime\prime}$,
kümeler$^{\prime\prime\prime\prime}$ vb. gerekir. Üs işaretlerinin
çoğalmasını önlemek için bunları kümeler$_0$, kümeler$_1$, kümeler$_2$,
kümeler$_3$, kümeler$_4$, \ldots{} olarak düşünmek daha kolaydır. Bu da bizi
(basit) tipler teorisine götürür.

Böyle bir teoriye karşı birkaç açık itiraz vardır (bunların \emph{ezici}
itirazlar olduğu açık olmasa da): Elde edilen teoriyi kullanmak zahmetlidir;
farklı türde çok sayıda nesne varsayar; ayrıca yüklemsel olmayan tanımların
sonuçta gerçekten \emph{o kadar kötü} olduğu açık değildir.

Son noktayı belirginleştirmek için \citeauthor{Ramsey1925}'nin şu
yorumunu düşünün:
\begin{quote}
  Bir adamdan bir grubun en uzunu diye söz edebilir, böylece hiçbir kısır
  döngü olmaksızın onu kendisinin de üyesi olduğu bir bütünlük aracılığıyla
  belirleyebiliriz. \citep{Ramsey1925}
\end{quote}
Ramsey'in vurgusu, ``gruptaki en uzun adam''ın yüklemsel olmayan bir tanım
\emph{olduğu}, fakat açıkça tümüyle sorunsuz olduğudur.

Buna, bu durumda en uzun kişiyi \emph{yüklemsel} yollarla seçebileceğimiz
yanıtı verilebilir. Örneğin söz konusu adamı parmağımızla gösterebiliriz.
O zaman yüklemsel olmayan tanımlara itiraz, yalnızca yüklemsel olmayan yolla
belirlenebilen varlıklarla sınırlandırılmalıdır. Fakat o durumda bile böyle
bir ``özsel yüklemsel olmama''nın neden kötü olduğu konusunda daha fazlasını
duymamız gerekir.\footnote{Daha fazlası için bkz. \citet{Linnebo2010}.}

Tanımlarımızın bir nesneyi \emph{yaratmak} için bize bir tür tarif vermesini
istiyorsak yüklemsel olmayan tanımların son derece kötü olduğu kabul
edilebilir. Böyle bir tanım verildiğinde gerçekten kısır döngüye düşülür:
Yüklemsel olmayan biçimde belirtilmiş nesneyi yaratmak için önce bütün
nesneleri (bu nesnenin kendisi dâhil) yaratmak gerekir; çünkü nesne bütün
nesneler aracılığıyla belirtilmiştir. Böylece nesneyi daha onu yaratmadan
önce yaratmak gerekir. Fakat bu da yalnızca kümeleri \emph{yarattığımız}
şeyler olarak düşünüyorsak ``özsel olarak yüklemsel olmayan'' kümelere karşı
ciddi bir itirazdır. Muhtemelen böyle düşünmüyoruz.

Bu nedenle---iyi ya da kötü---yaygınlaşan yaklaşım, yüklemsel olup olmama
konusunda katı bir çizgi izlemez. Bunun yerine bugün birikimli-yinelemeli
yaklaşım olarak görülen yolu izler. Sonunda bu, kümelerimizi ``evrelere''
ayırmamızı sağlar---yüklemsel yaklaşımın varlıkları kümeler$_0$,
kümeler$_1$, kümeler$_2$, \ldots{} olarak katmanlara ayırmasına \emph{biraz}
benzer---fakat aralarında bir tür farkı varsaymayız.

\end{document}
